% !TEX root = task-1-3-2.tex
% Editable task bank. Append new BankTask blocks at the end of the document.
\documentclass[12pt,a4paper]{article}
\usepackage{physics-textbook}
\usepackage{graphicx,needspace}
\usetikzlibrary{decorations.pathmorphing}
\newif\ifShowSolutions
\ifdefined\TutorEdition
  \ShowSolutionstrue
\else
  \ShowSolutionsfalse % Change to \ShowSolutionstrue to display solutions.
\fi
\newcounter{banktask}
\newsavebox{\BankTaskBox}
\newcommand{\BankTaskContent}[3]{%
  \noindent\textbf{\thebanktask.}\quad #2\par
  \ifShowSolutions
    \Needspace{50mm}% Keep the solution heading with its opening text or diagram.
    #3\par
    \ifdefined\TutorEdition
      \ifcsname TutorSource#1\endcsname
        \nopagebreak\smallskip{\small\csname TutorSource#1\endcsname\par}%
      \fi
    \fi
  \fi
}
\newcommand{\BankTask}[3]{%
  \refstepcounter{banktask}\par\medskip
  \sbox{\BankTaskBox}{\begin{minipage}{\linewidth}%
    \BankTaskContent{#1}{#2}{#3}\end{minipage}}%
  \ifdim\dimexpr\ht\BankTaskBox+\dp\BankTaskBox\relax>\dimexpr\textheight-20mm\relax
    \Needspace{35mm}\BankTaskContent{#1}{#2}{#3}%
  \else
    \Needspace{\dimexpr\ht\BankTaskBox+\dp\BankTaskBox\relax}%
    \noindent\usebox{\BankTaskBox}\par
  \fi
  \medskip
}
\begin{document}
\BookSubsection{1.3.2}{Закон движения}


% Source 3CF615: https://ege.fipi.ru/bank/questions.php?search=1&proj=BA1F39653304A5B041B656915DC36B38&qid=3CF615
\BankTask{3CF615}{
Координата $x$ тела меняется с течением времени $t$ согласно закону $x=23+5t-2t^2$, где все величины выражены в СИ. Определите проекцию $a_x$ ускорения этого тела, в $\text{м}/\text{с}^2$.
}{
\smallskip\noindent\textbf{Решение от GPT}\par
Сравним $x=23+5t-2t^2$ с $x=x_0+v_{0x}t+\dfrac{a_xt^2}{2}$: $\dfrac{a_x}{2}=-2$. Следовательно, $a_x=-4\,\text{м}/\text{с}^2$.
}

% Source 7CFB3E: https://ege.fipi.ru/bank/questions.php?search=1&proj=BA1F39653304A5B041B656915DC36B38&qid=7CFB3E
\BankTask{7CFB3E}{
Материальная точка движется прямолинейно
с постоянным ускорением вдоль оси $Ox$. График зависимости её координаты от времени $x=x(t)$ изображён на рисунке. Определите проекцию $a_{x}$ ускорения этого тела. Результат выразите в $\text{м}/\text{с}^2$.
\begin{center}\begin{tikzpicture}[x=1cm,y=1cm,>=stealth]
\draw[thin,black!65] (-2,0) grid (2,3);
\draw[->] (-2.6,0)--(3,0) node[right] {$t,\,\text{с}$};
\draw[->] (0,0)--(0,3.6) node[above] {$x,\,\text{м}$};
\foreach \x in {-2,-1,0,1,2} \node[below] at (\x,0) {$\x$};
\foreach \y in {1,2,3} \node[left] at (0,\y) {$\y$};
\draw[thick,dashed,domain=-2.45:0,samples=40] plot (\x,{.5*\x*\x});
\draw[very thick,domain=0:2.45,samples=40] plot (\x,{.5*\x*\x});
\end{tikzpicture}\end{center}
}{
\smallskip\noindent\textbf{Решение от GPT}\par
Из графика $x_0=0$, $v_{0x}=0$, $x(2)=2\,\text{м}$. По $x=a_xt^2/2$ получаем $a_x=2x/t^2=2\cdot2/4=1\,\text{м}/\text{с}^2$.
}

% Source B28492: https://ege.fipi.ru/bank/questions.php?search=1&proj=BA1F39653304A5B041B656915DC36B38&qid=B28492
\BankTask{B28492}{
Тело массой $200\,\text{г}$ движется вдоль оси $Ox$, при этом его координата изменяется во времени в соответствии с формулой $x(t)=10+5t-3t^2$ (все величины выражены в СИ).
\TaskMatchingColumns[.53\linewidth]{Физические величины}{
А) проекция $v_x(t)$ скорости тела\par
Б) проекция $F_x(t)$ равнодействующей сил, приложенных к телу
}{Формулы}{
1) $5-6t$\par
2) $-1{,}2$\par
3) $-3$\par
4) $10+5t$\par
}\begin{center}\begin{tabular}{|c|c|}\hline А & Б \\\hline\rule{0pt}{1.2em}\hspace{1cm} & \hspace{1cm} \\\hline\end{tabular}\end{center}
}{
\smallskip\noindent\textbf{Решение от GPT}\par
$v_x=dx/dt=5-6t$ (формула $1$), $a_x=-6\,\text{м}/\text{с}^2$. Тогда $F_x=ma_x=0{,}2\cdot(-6)=-1{,}2\,\text{Н}$ (формула $2$). Ответ: $12$.
}

% Source 69AD85: https://ege.fipi.ru/bank/questions.php?search=1&proj=BA1F39653304A5B041B656915DC36B38&qid=69AD85
\BankTask{69AD85}{
Выражения А и Б определяют зависимость координат двух тел от времени. Установите соответствие между зависимостью координаты тела от времени и зависимостью проекции скорости от времени для этого же тела. Все величины заданы в СИ.
\TaskMatchingColumns[.53\linewidth]{Координата}{
А) $x=10+5t+2t^2$\par
Б) $x=10+2t^2$
}{Скорость}{
1) $v_x=5+4t$\par
2) $v_x=5+2t$\par
3) $v_x=4t$\par
4) $v_x=2t$\par
}\begin{center}\begin{tabular}{|c|c|}\hline А & Б \\\hline\rule{0pt}{1.2em}\hspace{1cm} & \hspace{1cm} \\\hline\end{tabular}\end{center}
}{
\smallskip\noindent\textbf{Решение от GPT}\par
Скорость --- производная координаты: в случае А $v_x=5+4t$ (формула $1$), в случае Б $v_x=4t$ (формула $3$). Ответ: $13$.
}

% Source 35210: https://phys-ege.sdamgia.ru/problem?id=35210
\BankTask{35210}{
Бегун бежит прямолинейно с постоянной скоростью $4\,\text{м/с}$. В момент, когда бегун поравнялся с неподвижным велосипедистом, велосипедист начинает двигаться в ту же сторону равноускорено с ускорением $1\,\text{м/с}^{2}$. Найти скорость велосипедиста в момент, когда он догонит бегуна.
}{
\smallskip\noindent\textbf{Решение}\par
Будем считать началом отсчета точку, когда бегун поравнялся с неподвижным велосипедистом. Тогда уравнение движения бегуна при равномерном прямолинейном движении имеет вид:

\[x_1= v _1t.\]

Движение велосипедиста  — равноускоренное прямолинейное без начальной скорости  — описывается уравнением:

\[x_2= \frac{at ^2}{2} .\]

В момент встречи их координаты будут одинаковыми $v _1t= \frac{at ^2}{2} .$ Тогда велосипедист догонит бегуна в момент времени: 

\[t= \frac{2 v _1}{a} .\]

Скорость велосипедиста в этот момент времени будет равной:

\[v _2=at=2 v _1=2 \cdot 4=8 \,\text{м/с} .\]

Ответ: $8\,\text{м/с}$.
}

\end{document}
